\scanpage{61}
\mspage{30}
En fait, je conjecture que dans la conj. B, il est $\unsure{=|\overline m}$ vrai que $\Gamma^\circ$ agissant sur $X'=\widehat{\widetilde U}$ a un pt fixe et un seul (ce qui est plus fort, que ce point \inserted{fixe} sera nécessairement fixe par $\Gamma$). Ceci \deleted{signifie}\inserted{implique}, dans le cas des pts fixes à distance finie, qu'on a alors $\pi^{\Gamma^\circ}=(1)$, comme il se devrait en général --- et dans le cas du pt fixe à l'infini, que
\[
\pi^{\Gamma^\circ}\subset\operatorname{Norm}_{\pi}(L_i)=L_i
\]
donc la relation $\pi^{\unsure{\Gamma_0|\Gamma^\circ}}=L_i$ ``hautement plausible'' de tantôt !

Du coup je me demande d'ailleurs si ces énoncés ne devraient pas être valables sans partir d'une action de $\Gamma$ \deleted{\unsure{tt}} entier, mais seulement de $\Gamma^\circ$ (ou d'un sous-groupe contenant un sous-groupe ouvert de $\Gamma^\circ$, plus gén.) --- ce qui signifierait plus ou moins que
\margin{60}
